\section{AI/ML \& Forecasting Intelligence}
\label{sec:ai_ml_forecasting}

\subsection{Forecasting Objective}
\label{subsec:forecasting_objective}

The forecasting engine serves as the predictive foundation of GridSense AI, generating multi-step renewable power generation forecasts over 24-hour, 48-hour, and 72-hour planning horizons.

For each hourly time step $t \in \{1, 2, \dots, H\}$ across the horizon $H$, the engine estimates expected power generation $\hat{y}_t$ using:
\begin{itemize}[leftmargin=1.8em, itemsep=0.2em]
    \item \textbf{Historical generation patterns} capturing diurnal cycles, seasonal baselines, and auto-regressive momentum.
    \item \textbf{Meteorological forecast parameters} representing environmental drivers of solar irradiance and wind dynamics.
    \item \textbf{Site-level parameters} specifying installed physical capacity and conversion characteristics.
\end{itemize}

Rather than treating forecasting as an isolated mathematical regression task, the objective is to produce structured forecast trajectories and uncertainty envelopes that directly inform control-room dispatch decisions.

\subsection{Forecasting Approach}
\label{subsec:forecasting_approach}

For the hackathon prototype, GridSense AI adopts a machine learning framework based on gradient-boosted decision trees using \textbf{XGBoost} (eXtreme Gradient Boosting). XGBoost is well-suited for tabular renewable time-series data due to its computational efficiency, robust handling of non-linear feature interactions, and strong performance without requiring massive training clusters. Alternative sequence models (such as LSTM networks or Facebook Prophet) represent natural extensions for future comparative benchmarking.

The model learns complex non-linear mappings between renewable power output and multi-modal feature representations (\cref{tab:feature_engineering}).

\begin{table}[htbp]
    \centering
    \small
    \renewcommand{\arraystretch}{1.2}
    \caption{Feature Engineering Categories and Operational Purpose.}
    \label{tab:feature_engineering}
    \begin{tabularx}{\textwidth}{L{2.8cm} Y Y}
        \toprule
        \textbf{Input Category} & \textbf{Engineered Feature Set} & \textbf{Operational Purpose} \\
        \midrule
        \textbf{Weather Features} & Global Horizontal Irradiance (GHI), DNI, Cloud Cover \%, Temperature, Wind Speed ($v_{\text{hub}}$), Air Density & Quantify environmental energy flux and atmospheric attenuation. \\
        \textbf{Historical Features} & Lagged Generation ($y_{t-1}, y_{t-24}, y_{t-48}$), Rolling Means/Std, Hour of Day, Day of Year & Capture diurnal generation baselines, autocorrelation, and seasonality. \\
        \textbf{Site Features} & Nameplate Capacity ($P_{\text{rated}}$), Panel Tilt/Azimuth, Hub Height, Inverter Power Limits & Normalize generation curves against physical plant constraints. \\
        \bottomrule
    \end{tabularx}
\end{table}

\subsection{Forecast + Uncertainty}
\label{subsec:forecast_uncertainty}

A central premise of GridSense AI is that a deterministic point forecast is insufficient for reliable grid dispatch. In real-world power systems, identical point predictions can carry drastically different operating risks depending on meteorological stability.

Instead of reporting a solitary point value, GridSense AI generates an \textbf{uncertainty envelope}:
\begin{center}
    \texttt{Expected Generation } $= 80\text{ MW} \quad \big\vert \quad \texttt{Prediction Range } [p_{10}, p_{90}] = [72\text{ MW}, 88\text{ MW}] \quad \big\vert \quad \texttt{Confidence } = \text{Moderate}$
\end{center}

\begin{figure}[htbp]
    \centering
    \begin{tikzpicture}[xscale=0.82, yscale=0.68]
        % Background grid
        \draw[very thin, color=borderGray!50, step=1cm] (0,0) grid (12,5);
        \draw[->, thick, color=darkSlate] (0,0) -- (12.3,0) node[right] {\footnotesize Horizon (Hours)};
        \draw[->, thick, color=darkSlate] (0,0) -- (0,5.3) node[above] {\footnotesize Generation (MW)};
        
        % Ticks
        \foreach \x/\label in {0/00:00, 3/06:00, 6/12:00, 9/18:00, 12/24:00}
            \draw (\x,0.1) -- (\x,-0.1) node[below] {\scriptsize \label};
        \foreach \y/\label in {1/20, 2/40, 3/60, 4/80, 5/100}
            \draw (0.1,\y) -- (-0.1,\y) node[left] {\scriptsize \label};

        % Shaded Uncertainty Envelope (p10 to p90 band)
        \fill[accentTeal!18] 
            (0,0) -- 
            plot [smooth] coordinates {(2.5,0) (4,1.8) (6,4.8) (7.5,3.6) (9,1.6) (10,0)} --
            plot [smooth] coordinates {(10,0) (9,0.4) (7.5,1.8) (6,3.6) (4,0.6) (2.5,0)} -- cycle;

        % Net Demand Curve
        \draw[dashed, thick, color=darkSlate!70] plot [smooth] coordinates {(0,2.2) (3,1.8) (6,3.2) (9,4.5) (12,2.8)};
        \node[color=darkSlate!80, font=\scriptsize] at (10.4,4.7) {Scheduled Demand};

        % Expected Forecast Line (p50)
        \draw[line width=1.3pt, color=accentTeal] plot [smooth] coordinates {(0,0) (2.5,0) (4,1.2) (6,4.2) (7.5,2.7) (9,0.9) (10,0)};
        \node[color=accentTeal, font=\scriptsize\bfseries] at (6,4.5) {Expected Forecast ($p_{50}$)};

        % Highlighted Risk / Deficit Zone (Evening Peak)
        \draw[fill=accentRose!20, draw=accentRose, dashed, line width=0.8pt] (7.5,0.2) rectangle (9.8,4.6);
        \node[color=accentRose, font=\scriptsize\bfseries, align=center] at (8.65,2.5) {High Shortfall\\Risk Window};
    \end{tikzpicture}
    \caption{Conceptual 24-hour generation forecast with probabilistic prediction intervals ($p_{10}$--$p_{90}$) and highlighted shortfall risk window (illustrative).}
    \label{fig:forecast_uncertainty_chart}
\end{figure}

As shown in \cref{fig:forecast_uncertainty_chart}, the uncertainty envelope allows the decision layer to differentiate between stable meteorological regimes (tight prediction bands) and volatile atmospheric transitions (broad prediction intervals).

\subsection{Forecast Risk}
\label{subsec:forecast_risk}

Forecast uncertainty becomes operationally critical when potential deviations threaten grid balance. GridSense AI evaluates each planning interval using structured risk indicators:
\begin{itemize}[leftmargin=1.8em, itemsep=0.2em]
    \item \textbf{Expected Output Magnitude:} Anticipated generation relative to nominal capacity.
    \item \textbf{Dispersion / Volatility Index:} Width of the prediction interval ($p_{90} - p_{10}$), indicating meteorological uncertainty.
    \item \textbf{Surplus Imbalance Risk:} Probability of generation exceeding local hosting and export capacity ($P_{\text{gen}} > P_{\text{demand}} + P_{\text{export}}$).
    \item \textbf{Shortfall Imbalance Risk:} Probability of generation falling below critical commitment thresholds ($P_{\text{gen}} < P_{\text{demand}} - P_{\text{storage}}$).
    \item \textbf{Resource Adequacy Factor:} Comparison of potential shortfalls against online storage and flexible backup capacity.
\end{itemize}

Under this framework, a time window with moderate generation but severe forecast dispersion is flagged with higher operational priority than a calm, highly predictable low-generation window.

\subsection{Explainable Forecasts}
\label{subsec:explainable_forecasts}

To build operator trust and facilitate rapid situational awareness, the system translates feature dynamics and model outputs into clear, natural-language explanations:
\begin{center}
    \emph{``Solar generation is forecast to drop 42\% between 14:00 and 16:00 due to an incoming cloud front increasing cloud cover to 78\%. Shortfall risk is moderate; local battery storage is sufficient to bridge the gap.''}
\end{center}

This explainability layer connects raw predictive mechanics to operational reality without burdening control-room operators with model internal parameters.

\subsection{Model Evaluation}
\label{subsec:model_evaluation}

Rigorous offline validation is essential before deploying predictive models for decision support. GridSense AI structures model evaluation using standard time-series regression metrics:
\begin{itemize}[leftmargin=1.8em, itemsep=0.2em]
    \item \textbf{Mean Absolute Error (MAE):} Measures average absolute prediction error in megawatts:
    \begin{equation}
        \text{MAE} = \frac{1}{N} \sum_{i=1}^N |y_i - \hat{y}_i|
    \end{equation}
    \item \textbf{Root Mean Squared Error (RMSE):} Penalizes large, abrupt deviations heavily, reflecting operational costs associated with large unpredicted ramps:
    \begin{equation}
        \text{RMSE} = \sqrt{\frac{1}{N} \sum_{i=1}^N (y_i - \hat{y}_i)^2}
    \end{equation}
    \item \textbf{Mean Absolute Percentage Error (MAPE):} Evaluates normalized percentage error across non-zero daylight/generation periods:
    \begin{equation}
        \text{MAPE} = \frac{100\%}{N} \sum_{i=1}^N \left| \frac{y_i - \hat{y}_i}{y_i} \right| \quad (\text{for } y_i > \epsilon)
    \end{equation}
\end{itemize}

To prevent \textbf{data leakage}, model evaluation employs strict expanding-window or sliding-window time-aware cross-validation, guaranteeing that future meteorological data never informs past training instances. Concrete numerical benchmark results will be reported upon completion of model training.

\subsection{Forecasting Pipeline}
\label{subsec:forecasting_pipeline}

\Cref{fig:forecasting_pipeline} illustrates the structured, step-by-step pipeline connecting raw environmental feeds to operational risk indicators.

\begin{figure}[htbp]
    \centering
    \begin{tikzpicture}[
        node distance=0.38cm,
        pipeBox/.style={rectangle, rounded corners=3pt, draw=accentTeal!80, fill=boxBg, text width=7.6cm, align=left, font=\scriptsize, inner sep=3.5pt},
        arrow/.style={-{Stealth[scale=0.8]}, thick, color=darkSlate}
    ]
        \node[pipeBox] (p1) {\textbf{1. Multi-Modal Ingestion} \hfill \color{darkSlate!70}Historical Output, Weather, Site Params};
        \node[pipeBox, below=of p1] (p2) {\textbf{2. Data Preparation} \hfill \color{darkSlate!70}Imputation, Resampling, Outlier Removal};
        \node[pipeBox, below=of p2] (p3) {\textbf{3. Feature Engineering} \hfill \color{darkSlate!70}Lags, Rolling Stats, Solar Geometry, Ramps};
        \node[pipeBox, below=of p3] (p4) {\textbf{4. XGBoost Model Execution} \hfill \color{darkSlate!70}Multi-Step 24h--72h Horizon Regression};
        \node[pipeBox, below=of p4] (p5) {\textbf{5. Generation Forecast Output} \hfill \color{darkSlate!70}Hourly Point Trajectory ($\hat{y}_{t}$)};
        \node[pipeBox, below=of p5] (p6) {\textbf{6. Uncertainty Estimation} \hfill \color{darkSlate!70}Prediction Bands ($p_{10}, p_{50}, p_{90}$), Confidence};
        \node[pipeBox, below=of p6] (p7) {\textbf{7. Operational Risk Indicators} \hfill \color{darkSlate!70}Surplus/Shortfall Flags, Volatility Score};
        \node[pipeBox, below=of p7, draw=accentTeal, fill=accentTeal!15] (p8) {\textbf{8. Decision \& Optimization Engine} \hfill \color{darkSlate!70}\textbf{Storage Dispatch \& Grid Action}};

        \draw[arrow] (p1) -- (p2);
        \draw[arrow] (p2) -- (p3);
        \draw[arrow] (p3) -- (p4);
        \draw[arrow] (p4) -- (p5);
        \draw[arrow] (p5) -- (p6);
        \draw[arrow] (p6) -- (p7);
        \draw[arrow] (p7) -- (p8);
    \end{tikzpicture}
    \caption{Forecasting Intelligence Pipeline from multi-modal ingestion to the decision engine.}
    \label{fig:forecasting_pipeline}
\end{figure}

\subsection{Phase 4 Conclusion}
\label{subsec:phase4_conclusion}

The AI/ML layer establishes the predictive foundation of GridSense AI. 

Its primary innovation lies not simply in generating numerical time series, but in delivering forecasts coupled with uncertainty intervals, operational risk metrics, and natural-language explanations. This enables the downstream Decision and Optimization Engine (detailed in Phase 5) to balance risk and prescribe cost- and carbon-optimized grid interventions.
